\chapter{NUMA Architectures}
\label{chap:numa}

On a laptop with one CPU socket, almost all RAM is ``local.'' On multi-socket servers, 
each socket has its own memory controllers: accessing \highlight{local} RAM is faster than 
crossing to another socket's RAM. That design is \highlight{NUMA} 
(\emph{Non-Uniform Memory Access}). HFT boxes are often multi-socket.
\par\noindent\textbf{Interview tip:} keep \highlight{threads, memory, and NICs} on the same node when it matters.

\section{NUMA Fundamentals}
\label{sec:numa-fundamentals}
A \highlight{NUMA node} is typically one CPU socket plus the DRAM attached to it. 
Cores on that socket reach local DRAM with lower latency (and often higher bandwidth) 
than DRAM attached to another socket. Inter-socket links (e.g.\ Intel UPI / older QPI) 
carry remote traffic and add delay.

\subsection{Local Memory}
\label{subsec:numa-local-memory}
\highlight{Local memory} is RAM attached to the same node as the core executing the load/store. 
Prefer allocating hot structures (order books, ring buffers, decode scratch) on the node 
where the owning thread runs.
\newline
\textbf{Noise removed:} the inter-socket interconnect round-trip on every hot load/store. 
Local DRAM is still hundreds of cycles versus cache, but it is the best DRAM you can get 
on that box --- and it keeps bandwidth on the local memory controllers.

\subsection{Remote Memory}
\label{subsec:numa-remote-memory}
\highlight{Remote memory} is RAM on another node. Every remote access pays the interconnect 
tax and can increase jitter. Cross-node false sharing / coherence traffic is especially costly 
(Chapter~\ref{chap:cacheperformance}).
\newline
\textbf{Noise removed} by staying local: UPI/QPI (or equivalent) hops and remote memory-controller 
contention that show up as fat tails even when median latency looks fine. Cost of fixing it: 
explicit binding and allocation discipline --- the OS will not guess your topology for you.
\par\noindent\textbf{Interview tip:} local = same node as the core; remote = another node's DRAM --- 
name both latency and the coherence tax when lines bounce across sockets.

\section{NUMA Performance}
\label{sec:numa-performance}

\subsection{Memory Access Latency}
\label{subsec:numa-memory-latency}
\par\noindent\textbf{Interview tip:} remote DRAM can be on the order of 
\highlight{$\sim$1.5$\times$--2$\times+$} local latency (platform-dependent; measure on the box). 
Bandwidth can suffer too when many cores hammer the interconnect.
\newline
Classic failure mode: thread pinned on node 0, but \texttt{malloc}/\texttt{new} first-touched 
pages on node 1 (or the OS placed them poorly) $\Rightarrow$ every hot access is remote.
\newline
That split is easy to miss in code review: affinity looks correct in \texttt{taskset}, yet 
profiles show remote DRAM hits because allocation happened on a different node (or before 
pinning). Another variant: NIC on node 0, strategy on node 1 --- packets cross the 
interconnect twice (DMA placement + CPU access) unless buffers and threads share a node.
\newline
\textbf{Noise removed} by first-touch / \texttt{membind} after pinning on the NIC's node: 
systematic remote DRAM on the hot path. See Section~\ref{subsec:numa-allocation-locality}.

\section{NUMA-Aware Programming}
\label{sec:numa-aware-programming}
Goals:
\begin{enumerate}
    \item Run the thread on the right node (CPU binding).
    \item Allocate its hot memory on that same node (allocation locality).
    \item Prefer the NIC that sits on that node (PCI proximity; Chapter~\ref{chap:linuxprocesses} topology notes).
\end{enumerate}

\subsection{Allocation Locality}
\label{subsec:numa-allocation-locality}
Linux often uses \highlight{first-touch}: the node of the CPU that first writes a page 
owns that page. Patterns:
\begin{itemize}
    \item Pin the thread first, then allocate / initialise large buffers on that thread.
    \item Or allocate explicitly with \texttt{libnuma} / \texttt{numactl}.
\end{itemize}
Command-line control (process-wide):
\begin{lstlisting}[language=bash]
# Run on CPUs of node 0; prefer allocating on node 0
numactl --cpunodebind=0 --membind=0 ./trading_engine

# Show topology
numactl --hardware
lscpu | grep NUMA
\end{lstlisting}
Library sketch (link with \texttt{-lnuma}; package \texttt{libnuma-dev} on Debian/Ubuntu):
\begin{lstlisting}[language=C++]
#include <numa.h>

void* alloc_on_node(std::size_t bytes, int node)
{
    if (numa_available() < 0) {
        return nullptr;  // not a NUMA system / lib unavailable
    }
    return numa_alloc_onnode(bytes, node);
}

void free_numa(void* p, std::size_t bytes)
{
    numa_free(p, bytes);
}
\end{lstlisting}
On a \highlight{single-node} machine (common laptops), \texttt{numactl --hardware} shows one node; 
NUMA tuning still matters less than pinning and IRQ isolation, but the vocabulary remains 
interview-relevant for servers.

\subsection{NUMA Thread and CPU Binding}
\label{subsec:numa-thread-cpu-binding}
For core-level pinning with \texttt{sched\char`_setaffinity()} / \texttt{pthread\char`_setaffinity\char`_np}, 
see Section~\ref{subsec:cpu-affinity-thread-pinning}. 
This section covers placing threads on CPUs \highlight{within a NUMA node} so they stay next 
to their data (and often next to the NIC).
\newline
Pinning a thread without placing its memory (and ideally its NIC) on the same node only 
removes migration noise --- remote DRAM remains. The workflow below closes that gap.
\newline
Workflow:
\begin{enumerate}
    \item Find the NIC's node: \texttt{/sys/bus/pci/devices/<addr>/numa\char`_node} 
          (\texttt{-1} often means ``unspecified,'' treat as node 0 on single-node hosts).
    \item List CPUs on that node: \texttt{/sys/devices/system/node/nodeN/cpulist}.
    \item Pin trading threads to those CPUs (avoid core 0 when IRQs land there).
    \item Allocate hot memory on the same node (\texttt{membind} / first-touch / \texttt{numa\char`_alloc\char`_onnode}).
\end{enumerate}
\textbf{Noise removed:} remote memory accesses and cross-node packet paths that survive 
``we pinned the thread'' alone. Cost: less OS freedom to balance load; you own placement.
\newline
\texttt{libnuma} helpers such as \texttt{numa\char`_run\char`_on\char`_node(node)} restrict the calling 
thread to that node's CPUs; combine with explicit affinity when you need a \highlight{specific} core.
\par\noindent\textbf{Interview tip:} \highlight{pin compute next to memory next to the NIC} --- same NUMA node --- 
unless you have measured that another layout wins.
